% ============================================================
%  UPC 论文模板 · 官方 elsarticle 款
%  来源：官方 template.tex（Elsevier elsarticle），清理了占位内容，
%        按金牌论文惯例补齐注释与结构说明。
%  适用：想要和官方示例完全同源、依赖 TeX Live/MiKTeX 自带 elsarticle.cls。
%  编译：pdflatex（两遍，含 bibtex 则三遍）
%  搭配：总览页「论文模板工坊」对比选择；写作流程见《一篇金牌论文的完整历程》
% ============================================================
\documentclass[final,3p,times,onecolumn]{elsarticle}
% 其他版式开关（官方注释保留）：
%   \documentclass[preprint,12pt]{elsarticle}      草稿大字距
%   \documentclass[final,5p,times,twocolumn]{elsarticle}  期刊双栏

\usepackage[fleqn]{amsmath}
\usepackage{amssymb}
\usepackage{graphicx}
\usepackage{float}
\usepackage{booktabs}                 % 三线表：\toprule\midrule\bottomrule
\usepackage{caption}
\captionsetup{font=small,labelfont=bf}
\usepackage{mhchem}                   % 化学式（如需）
\usepackage[usenames,dvipsnames]{xcolor}
\usepackage[colorlinks=true,linkcolor=blue!60!black,citecolor=blue!60!black,urlcolor=blue!60!black]{hyperref}

% 金牌高频宏
\newcommand{\dd}{\mathrm{d}}
\newcommand{\vect}[1]{\boldsymbol{#1}}
\newcommand{\unit}[1]{\,\mathrm{#1}}

\biboptions{numbers,sort&compress}

\begin{document}

\begin{frontmatter}

% 匿名规范：只写队号与题号，不出现姓名/学校
\title{Your Descriptive Paper Title Here}
\author{Team \#1234}
\address{Problem A}
\date{\today}

% 摘要一页：符号定义 → 问题两难 → 方法路线 → 带数字的结论（评委只读摘要也能拿结果）
\begin{abstract}
We define the key quantities ($R \approx 2.5\times10^{4}\unit{km}$,
$v_0 = 30\unit{km/s}$) and identify the central trade-off: if [A] is too large,
[failure 1]; if too small, [failure 2].
Section~1 establishes a baseline costing 95\% of initial mass.
Section~2 develops the model with data-anchored parameters, integrated by
[method] and verified by step-halving to 4--5 significant figures.
The admissible window is
\[
  \boxed{\;8.8296^\circ < \theta < 8.8549^\circ\;}
\]
We therefore \emph{advise against} the maneuver and identify [parameter] as the
most promising lever.
\end{abstract}

% \begin{keyword}
%   keyword1 \sep keyword2
% \end{keyword}

\end{frontmatter}

% ---------- 1 基线：先算"什么都不做"的成本，给后文量化参照 ----------
\section{Baseline: The Cost of Doing Nothing}
We first compute the cost of [baseline option]:
\begin{equation}
  E = \tfrac12 m v_0^2 - \frac{GMm}{R},
\end{equation}
which is positive --- the unmodified trajectory is hyperbolic. A circular orbit
requires $\Delta v = 13.62\unit{km/s}$, i.e.\ by the rocket equation
\begin{equation}
  m_f = m_0\, e^{-\Delta v / v_{ex}} \approx 121\unit{kg},
  \label{eq:rocket}
\end{equation}
a 95.15\% mass cost. This anchors all comparisons below.

% ---------- 2 模型：理想 → 关键修正 → 扩展，每层给量级论证 ----------
\section{Model}
\subsection{Governing Equations}
Quadratic drag and Newtonian gravity give
\begin{equation}
  m\ddot{\vect{r}} = -\frac{GMm}{r^3}\vect{r}
  \;-\; \sigma\,\frac{P(h)}{T(h)}\,|\dot{\vect r}|\,\dot{\vect{r}},
  \qquad \sigma \equiv \frac{c_d \mu A}{2k_B},
  \label{eq:eom}
\end{equation}
where $P(h)$ and $T(h)$ are anchored at two literature points each~\cite{Key1993}.

\subsection{Numerical Method}
Equation~\eqref{eq:eom} is integrated by [method] with $\Delta t = 10^{-3}\unit{s}$;
the run terminates at [condition]. Full code in Appendix~B.

\subsection{Operationalizing the Criterion}
A ``significant'' saving is not quantitative; we require $m_f \ge 0.8\,m_0$,
which bounds the admissible [quantity] to
\begin{equation}
  q_{min} \le q \le q_{max}.
\end{equation}

% ---------- 3 结果：单调性 + 反例夹逼 ----------
\section{Results}
Exit speed increases monotonically with [parameter] (App.~A), so the admissible
set is one interval. Bracketing by counterexample simulations:
\begin{equation}
  8.8295^\circ < q_{min} < 8.8296^\circ .
\end{equation}

\begin{figure}[htbp]
  \centering
  \includegraphics[width=0.85\linewidth]{Figure1}
  \caption{[This figure proves one point: ...]}
  \label{fig:1}
\end{figure}

\begin{table}[htbp]
  \caption{[One claim per table.]}
  \label{tab:1}
  \begin{tabular}{lcccc}
    \toprule
    phase & 1 & 2 & 3 & 4 \\
    \midrule
    gold   & $-$   & $-0.4$ & $-0.4$  & $-1.9$ \\
    silver & $-$   & $-0.8$ & $1.9$   & $-1.1$ \\
    \bottomrule
  \end{tabular}
\end{table}

% ---------- 4 误差三件套：收敛性/极限情形/灵敏度 ----------
\section{Uncertainty}
Halving the step changes the result in the fifth digit; we claim 4--5
significant figures, worst case 4. Unmodeled: [rotation, winds, shape].

% ---------- 5 结论：区间 + 与基准对比 + 敢写负结论 ----------
\section{Conclusion}
The admissible window is $0.0253^\circ$ wide --- a 75\% fuel saving at the
price of an impractically tight corridor. We advise against the maneuver;
[targeting X] would relax the window considerably.

\appendix
\section{Derivation of the Rocket Equation}
\section{Simulation Code}

% ---------- 参考文献：金牌普遍 5–15 条，网络来源注明访问日期 ----------
\bibliographystyle{elsarticle-num}
\bibliography{ref}
% 若不用 BibTeX，可改用 \begin{thebibliography}{9} 手排

\end{document}
